
\documentclass{article}
\usepackage{colortbl}
\usepackage{makecell}
\usepackage{multirow}
\usepackage{supertabular}

\begin{document}

\newcounter{utterance}

\centering \large Interaction Transcript for game `bbh', experiment `date\_understanding', episode 107 with reflexion{-}llama{-}3.1{-}8b{-}instruct.
\vspace{24pt}

{ \footnotesize  \setcounter{utterance}{1}
\setlength{\tabcolsep}{0pt}
\begin{supertabular}{c@{$\;$}|p{.15\linewidth}@{}p{.15\linewidth}p{.15\linewidth}p{.15\linewidth}p{.15\linewidth}p{.15\linewidth}}
    \# & \multicolumn{2}{c}{Player} && \multicolumn{2}{c}{Game Master} \\
    \hline

    \theutterance \stepcounter{utterance}  
    & & & \multicolumn{4}{p{0.6\linewidth}}{
        \cellcolor[rgb]{0.9,0.9,0.9}{
            \makecell[{{p{\linewidth}}}]{
                \texttt{\tiny{[P1$\langle$GM]}}
                \texttt{You have attempted this same task before and did not succeed. These are the plans you wrote after those attempts.} \\
\texttt{Trial \#1: New plan: I failed to consider the fact that the date one week from today in MM/DD/YYYY is simply the next date in the sequence, without necessarily requiring any complex calculations. Next time, I will take a more straightforward approach by simply adding one week to the given date, without overcomplicating the problem.} \\
\\ 
\texttt{Trial \#2: New plan: I failed to consider the fact that the date one week from today in MM/DD/YYYY is simply the next date in the sequence, without necessarily requiring any complex calculations. However, this time I also failed to apply this straightforward approach correctly. Next time, I will take a more careful and systematic approach by adding one week to the given date, considering the month and year transitions, and making sure to account for the correct number of days in each month.} \\
\\ 
\texttt{Trial \#3: New plan: I failed to consider the fact that the date one week from today in MM/DD/YYYY is simply the next date in the sequence, without necessarily requiring any complex calculations. However, this time I also failed to apply this straightforward approach correctly. Next time, I will take a more careful and systematic approach by adding one week to the given date, considering the month and year transitions, and making sure to account for the correct number of days in each month. } \\
\\ 
\texttt{To achieve this, I will break down the process into smaller steps:} \\
\\ 
\texttt{1. Identify the current date and its components (month, day, year).} \\
\texttt{2. Determine the number of days in the current month.} \\
\texttt{3. Add 7 days to the current date, considering the month and year transitions.} \\
\texttt{4. If the resulting date falls in the next month, adjust the month and year accordingly.} \\
\texttt{5. If the resulting date falls in the next year, adjust the year and month accordingly.} \\
\\ 
\texttt{By following these steps, I will ensure that I accurately calculate the date one week from the given date, without overcomplicating the problem.} \\
\\ 
\\ 
\texttt{Infer the date from context.} \\
\\ 
\texttt{Q: Today is Christmas Eve of 1937. What is the date 10 days ago in MM/DD/YYYY?} \\
\texttt{Options:} \\
\texttt{(A) 12/14/2026} \\
\texttt{(B) 12/14/1950} \\
\texttt{(C) 12/14/2007} \\
\texttt{(D) 12/14/1937} \\
\texttt{(E) 07/14/1938} \\
\texttt{(F) 12/14/1988} \\
\texttt{A: (D)} \\
\\ 
\texttt{Q: Tomorrow is 11/12/2019. What is the date one year ago from today in MM/DD/YYYY?} \\
\texttt{Options:} \\
\texttt{(A) 09/04/2018} \\
\texttt{(B) 11/11/2018} \\
\texttt{(C) 08/25/2018} \\
\texttt{(D) 11/02/2018} \\
\texttt{(E) 11/04/2018} \\
\texttt{A: (B)} \\
\\ 
\texttt{Q: Jane and John married on Jan 2, 1958. It is their 5{-}year anniversary today. What is the date tomorrow in MM/DD/YYYY?} \\
\texttt{Options:} \\
\texttt{(A) 01/11/1961} \\
\texttt{(B) 01/03/1963} \\
\texttt{(C) 01/18/1961} \\
\texttt{(D) 10/14/1960} \\
\texttt{(E) 01/03/1982} \\
\texttt{(F) 12/03/1960} \\
\texttt{A: (B)} \\
\\ 
\texttt{Q: Today is Sep 9, 1909. What is the date one week from today in MM/DD/YYYY?} \\
\texttt{Options:} \\
\texttt{(A) 09/16/1952} \\
\texttt{(B) 09/16/1909} \\
\texttt{(C) 09/16/1876} \\
\texttt{(D) 08/22/1909} \\
\texttt{(E) 12/16/1909} \\
\texttt{(F) 10/14/1909} \\
\texttt{A:} \\
            }
        }
    }
     \\ \\

    \theutterance \stepcounter{utterance}  
    & \multicolumn{4}{p{0.6\linewidth}}{
        \cellcolor[rgb]{0.561,0.737,0.561}{
            \makecell[{{p{\linewidth}}}]{
                \texttt{\tiny{[P1$\rangle$GM]}}
                \texttt{Based on the context, the question is asking for the date one week from September 9, 1909. } \\
\\ 
\texttt{To find the date one week from September 9, 1909, we can simply add 7 days to the given date. } \\
\\ 
\texttt{September 9, 1909 + 7 days = September 16, 1909.} \\
\\ 
\texttt{So, the correct answer is:} \\
\\ 
\texttt{(B) 09/16/1909} \\
            }
        }
    }
    & & \\ \\

    \theutterance \stepcounter{utterance}  
    & & & \multicolumn{2}{p{0.3\linewidth}}{
        \cellcolor[rgb]{0.9,0.9,0.9}{
            \makecell[{{p{\linewidth}}}]{
                \texttt{\tiny{[GM$|$GM]}}
                \texttt{Target: (B)} \\
            }
        }
    }
    & & \\ \\

    \theutterance \stepcounter{utterance}  
    & & & \multicolumn{2}{p{0.3\linewidth}}{
        \cellcolor[rgb]{0.9,0.9,0.9}{
            \makecell[{{p{\linewidth}}}]{
                \texttt{\tiny{[GM$|$GM]}}
                \texttt{game\_result = LOSE} \\
            }
        }
    }
    & & \\ \\

    \theutterance \stepcounter{utterance}  
    & & \multicolumn{4}{p{0.6\linewidth}}{
        \cellcolor[rgb]{0.95,0.95,0.95}{
            \makecell[{{p{\linewidth}}}]{
                \texttt{\tiny{[P1$\leftrightarrow$P]}}
                \texttt{Post{-}episode Reflexion (not sent to GM):} \\
\texttt{New plan: I failed to apply the straightforward approach correctly in my previous attempts. This time, I will take a more careful and systematic approach by adding one week to the given date, considering the month and year transitions, and making sure to account for the correct number of days in each month. } \\
\\ 
\texttt{To achieve this, I will break down the process into smaller steps:} \\
\\ 
\texttt{1. Identify the current date and its components (month, day, year).} \\
\texttt{2. Determine the number of days in the current month.} \\
\texttt{3. Add 7 days to the current date, considering the month and year transitions.} \\
\texttt{4. If the resulting date falls in the next month, adjust the month and year accordingly.} \\
\texttt{5. If the resulting date falls in the next year, adjust the year and month accordingly.} \\
\\ 
\texttt{However, I also noticed that my previous attempts were too focused on the general approach, without considering the specific details of the given date. This time, I will make sure to take into account the specific month and year of the given date, and adjust my calculations accordingly.} \\
\\ 
\texttt{For example, if the given date is September 9, 1909, I will make sure to consider the fact that September has 30 days, and add 7 days to the current date, resulting in September 16, 1909. If the given date is a month with fewer than 30 days, I will make sure to adjust the calculation accordingly.} \\
\\ 
\texttt{By following these steps and considering the specific details of the given date, I will ensure that I accurately calculate the date one week from the given date, without overcomplicating the problem.} \\
            }
        }
    }
    & \\ \\

\end{supertabular}
}

\end{document}
